\documentclass[11pt,a4paper]{article}

\usepackage[T1]{fontenc}
\usepackage[utf8]{inputenc}
\usepackage{mathptmx}
\usepackage[scaled=.92]{helvet}
\usepackage{courier}
\usepackage{microtype}
\DisableLigatures[-]{encoding=*,family=tt*}  % keep -- as two hyphens in monospace text
\usepackage[margin=2.4cm]{geometry}
\usepackage{xcolor}
\usepackage{booktabs}
\usepackage{longtable}
\usepackage{array}
\usepackage{enumitem}
\usepackage{fancyhdr}
\usepackage{listings}
\usepackage[hidelinks]{hyperref}

\definecolor{accent}{HTML}{1F5C8B}
\definecolor{codebg}{HTML}{F5F7FA}
\definecolor{codeframe}{HTML}{C9D3DE}
\definecolor{comment}{HTML}{5C6B7A}

\hypersetup{pdftitle={sysinfo User Manual},pdfauthor={sysinfo},colorlinks=true,linkcolor=accent,urlcolor=accent}

\setlist{itemsep=2pt,topsep=4pt}
\setlength{\parindent}{0pt}
\setlength{\parskip}{6pt}
\renewcommand{\arraystretch}{1.2}

\lstdefinestyle{base}{
  basicstyle=\ttfamily\footnotesize,
  backgroundcolor=\color{codebg},
  frame=single, rulecolor=\color{codeframe},
  breaklines=true, breakatwhitespace=false,
  postbreak=\mbox{\textcolor{comment}{$\hookrightarrow$}\space},
  columns=fullflexible, keepspaces=true,
  aboveskip=8pt, belowskip=6pt,
  xleftmargin=4pt, xrightmargin=4pt, framexleftmargin=4pt,
  showstringspaces=false, upquote=true,
}
\lstdefinestyle{shell}{style=base,language=bash,commentstyle=\color{comment}\itshape,
  morekeywords={sysinfo,cmake,sudo,apt,ctest,jq,grep,diff}}
\lstdefinestyle{out}{style=base}
\lstdefinestyle{cpp}{style=base,language=C++,keywordstyle=\color{accent}\bfseries,commentstyle=\color{comment}\itshape}
\lstdefinestyle{cmake}{style=base,commentstyle=\color{comment}\itshape,
  morekeywords={add_subdirectory,target_link_libraries,FetchContent_Declare,FetchContent_MakeAvailable,find_package},
  keywordstyle=\color{accent}\bfseries}
\lstset{style=shell}

\newcommand{\cmd}[1]{\texttt{#1}}
\newcommand{\opt}[1]{\texttt{\detokenize{#1}}}
\newcommand{\key}[1]{\texttt{\detokenize{sys_#1}}}
\newcommand{\prog}{\texttt{sysinfo}}
\newcommand{\sysinfo}{\texttt{sysinfo}}

\pagestyle{fancy}
\fancyhf{}
\fancyhead[L]{\small\textcolor{comment}{sysinfo User Manual}}
\fancyhead[R]{\small\textcolor{comment}{version 1.0.0}}
\fancyfoot[C]{\small\thepage}
\renewcommand{\headrulewidth}{0.3pt}

\begin{document}

% ---------------------------------------------------------------------------
\begin{titlepage}
\vspace*{3.5cm}
{\Huge\bfseries\color{accent} sysinfo}\par
\vspace{4pt}
{\Large Print the specifications of the machine you are on}\par
\vspace{1.2cm}
{\large User Manual \quad$\cdot$\quad Version 1.0.0}\par
\vspace{2cm}
\begin{minipage}{0.8\textwidth}
\prog{} is a small command line program for Linux that reports the operating system, processor,
memory, storage, graphics, network, sensors, user environment and the toolchain it was built with.
It reads only \cmd{/proc}, \cmd{/sys} and a few system calls: no external programs, no root rights, no
network access. The output is available as readable text, JSON or \texttt{key=value} lines.
\end{minipage}
\vfill
{\small Primary platform: Ubuntu 24.04 LTS and newer \quad $\cdot$ \quad C++20 \quad $\cdot$ \quad CMake $\geq$ 3.20}
\end{titlepage}

\tableofcontents
\newpage

% ---------------------------------------------------------------------------
\section{Quick start}

\begin{lstlisting}
sudo apt install build-essential cmake      # once
cmake -S . -B build                         # configure (Release by default)
cmake --build build -j                      # compile
./build/sysinfo                             # print everything
./build/sysinfo -s cpu,memory               # only some sections
./build/sysinfo -f json -o machine.json     # machine readable, to a file
\end{lstlisting}

That is all that is needed. The rest of this manual describes the options (Section~\ref{sec:cli}), the tests (Section~\ref{sec:tests}), the
information that is reported (Section~\ref{sec:sections}), how to embed the collector in your own
program (Section~\ref{sec:library}) and what to do when something looks wrong
(Section~\ref{sec:trouble}).

% ---------------------------------------------------------------------------
\section{Installation}

\subsection{Requirements}

\begin{itemize}
  \item A C++20 compiler: GCC 11 or newer, or Clang 14 or newer with libstdc++. Ubuntu 24.04 ships GCC 13.
  \item CMake 3.20 or newer (Ubuntu 24.04 ships 3.28).
  \item For this manual only: a \LaTeX{} installation with \cmd{latexmk}
        (\cmd{sudo apt install latexmk texlive-latex-recommended texlive-latex-extra}).
\end{itemize}

No third-party libraries are used.

\subsection{Building}

\begin{lstlisting}
cmake -S . -B build
cmake --build build -j
ctest --test-dir build            # optional: 12 tests
\end{lstlisting}

If no build type is given, CMake configures a \texttt{Release} build. The build type and the compiler flags
are recorded in the executable and reported in the \texttt{build} section.

CMake options:

\begin{center}
\begin{tabular}{@{}lll@{}}
\toprule
Option & Default & Meaning \\
\midrule
\opt{SYSINFO_BUILD_CLI}          & ON  & build the \prog{} executable \\
\opt{SYSINFO_BUILD_TESTS}        & ON  & register the tests for \cmd{ctest} (top-level builds only) \\
\opt{SYSINFO_WARNINGS_AS_ERRORS} & OFF & compile with \cmd{-Werror} (\cmd{/WX} on MSVC) \\
\bottomrule
\end{tabular}
\end{center}

Example with another compiler and a debug build:

\begin{lstlisting}
cmake -S . -B build-clang -DCMAKE_CXX_COMPILER=clang++ -DCMAKE_BUILD_TYPE=Debug
\end{lstlisting}

\subsection{Installing}

\begin{lstlisting}
sudo cmake --install build                          # to /usr/local/bin
cmake --install build --prefix "$HOME/.local"       # without root
\end{lstlisting}

\subsection{Building this manual}

\begin{lstlisting}
cmake --build build --target manual                 # -> build/doc/sysinfo-manual.pdf
\end{lstlisting}

The target exists when \cmd{latexmk} is found at configure time. The example outputs in this manual are the
files in \cmd{doc/examples/}, generated by the program itself.

% ---------------------------------------------------------------------------
\section{Tests}
\label{sec:tests}

The project ships 12 command line tests, registered with CTest in \cmd{CMakeLists.txt}. They run the compiled \prog{}
executable and check its exit status and, for four of them, a pattern in its output. They are smoke and behaviour
tests of the command line tool: they confirm that every mode runs on the current machine and that wrong input is
rejected. They do not compare hardware values, so they pass on any computer.

\subsection{When the tests are available}

The tests are registered only when \emph{all} of the following hold:

\begin{itemize}
  \item \sysinfo{} is the top-level CMake project (they are not registered when it is used from another project through \texttt{add\_}\allowbreak\texttt{subdirectory});
  \item \opt{SYSINFO_BUILD_CLI} is \texttt{ON} (the tests need the executable);
  \item \opt{SYSINFO_BUILD_TESTS} is \texttt{ON} (the default).
\end{itemize}

To switch them off: \cmd{cmake -S . -B build -DSYSINFO\_BUILD\_TESTS=OFF}.

\subsection{Running the tests}

Build first, then run CTest on the build directory:

\begin{lstlisting}
cmake -S . -B build
cmake --build build -j
ctest --test-dir build                       # run all tests
\end{lstlisting}

A successful run ends with:

\begin{lstlisting}[style=out]
100% tests passed, 0 tests failed out of 12
\end{lstlisting}

The tests use the executable of the build directory; after changing the sources, rebuild before running them again.
Useful variants:

\begin{center}
\begin{tabular}{@{}p{7.4cm}p{7.6cm}@{}}
\toprule
Command & Effect \\
\midrule
\cmd{ctest --test-dir build -N}                     & list the tests without running them \\
\cmd{ctest --test-dir build --output-on-failure}    & print the program output of failed tests \\
\cmd{ctest --test-dir build -V}                     & verbose: show each command line and its output \\
\cmd{ctest --test-dir build -R json}                & only tests whose name matches the regular expression \\
\cmd{ctest --test-dir build -E bad}                 & all tests except those matching \\
\cmd{ctest --test-dir build -j4}                    & run four tests in parallel \\
\cmd{ctest --test-dir build --rerun-failed}         & repeat only the tests that failed last time \\
\bottomrule
\end{tabular}
\end{center}

With a multi-configuration generator (Visual Studio, Xcode) add the configuration: \cmd{ctest --test-dir build -C Release}.

\subsection{The tests}

\emph{Exit 0} means the test passes when \prog{} exits with status 0. \emph{Must fail} means the test passes when
the exit status is not 0 (CTest property \texttt{WILL\_FAIL}): these check that errors are detected. A \emph{pattern} is a
regular expression that must occur in the standard output (\texttt{PASS\_REGULAR\_EXPRESSION}); a test with a pattern passes
only if the pattern matches, whatever the exit status.

{\small
\begin{longtable}{@{}l p{4.9cm} p{7.2cm}@{}}
\toprule
Name & Command & Passes when / what it checks \\
\midrule \endhead
\texttt{help}     & \cmd{sysinfo --help}    & exit 0. The usage text is printed. \\
\texttt{version}  & \cmd{sysinfo --version} & output matches \texttt{\^{}sysinfo <version>}, where \texttt{<version>} is the one given to
                    \cmd{project()} in \cmd{CMakeLists.txt}. \\
\texttt{list}     & \cmd{sysinfo --list}    & exit 0. The section list is printed. \\
\texttt{text}     & \cmd{sysinfo}           & output matches\newline\texttt{System.*Processor.*Memory}: all sections are collected and printed in text
                    format, in display order. \\
\texttt{verbose}  & \cmd{sysinfo --verbose} & exit 0. The long-tail collection (all flags, pseudo file systems, virtual interfaces) runs. \\
\texttt{json}     & \cmd{sysinfo --format json} & output matches\newline\texttt{"sys\_os": } (key, colon, space); JSON is produced and contains the operating system key. \\
\texttt{kv}       & \cmd{sysinfo -f kv -s core} & output matches\newline\texttt{sys\_cpu\_logical\_cores=[0-9]+}: key/value output with a numeric core
                    count; selection by the \texttt{core} group. \\
\texttt{redacted} & \cmd{sysinfo -r -x build}   & exit 0. Redaction combined with a section exclusion. \\
\texttt{bad\_option}  & \cmd{sysinfo --no-such-option} & must fail. An unknown option is rejected. \\
\texttt{bad\_section} & \cmd{sysinfo -s nonsense}      & must fail. An unknown section name is rejected. \\
\texttt{bad\_format}  & \cmd{sysinfo -f xml}           & must fail. An unknown output format is rejected. \\
\texttt{unwritable\_file} & \texttt{sysinfo -o} \path{/nonexistent-dir/out.txt} & must fail. An output file that cannot be created is
                    reported as an error. \\
\bottomrule
\end{longtable}
}

\subsection{Limits of the tests}

\begin{itemize}
  \item The \emph{must fail} tests accept any non-zero status, so they do not distinguish the documented usage error (2) from the I/O
        error (1). To verify the exact code run the command by hand and read \cmd{echo \$?}; see Section~\ref{sec:cli}.
  \item The pattern tests inspect only the standard output. A pattern is a regular expression in CMake syntax, so changing the
        output wording (for example a section title) requires updating the pattern in \cmd{CMakeLists.txt}.
  \item The tests do not verify that the reported values are correct for the hardware. Compare the output with
        \cmd{lscpu}, \cmd{free}, \cmd{lsblk} or \cmd{ip addr} for that.
  \item No test checks colour output, because the output is not a terminal under CTest.
\end{itemize}

\subsection{Running the tests with sanitizers}

A debug build with the compiler sanitizers is a good way to find memory errors and undefined behaviour in the file parsing code:

\begin{lstlisting}
cmake -S . -B build-asan -DCMAKE_BUILD_TYPE=Debug -DSYSINFO_WARNINGS_AS_ERRORS=ON \
      -DCMAKE_CXX_FLAGS="-fsanitize=address,undefined -fno-sanitize-recover=all"
cmake --build build-asan -j
ctest --test-dir build-asan --output-on-failure
\end{lstlisting}

\subsection{Adding a test}

Add a line to the test block of \cmd{CMakeLists.txt}, for example a test that the \texttt{gpu} section can be printed alone and mentions its title:

\begin{lstlisting}[style=cmake]
add_test(NAME gpu_only COMMAND sysinfo -s gpu)
set_tests_properties(gpu_only PROPERTIES PASS_REGULAR_EXPRESSION "Graphics")
\end{lstlisting}

Re-run CMake (it does so automatically on the next build) and check that the new test is listed with \cmd{ctest --test-dir build -N}.

\subsection{Tests and continuous integration}

CTest returns a non-zero exit status when any test fails, so the single command \cmd{ctest --test-dir build --output-on-failure} is
enough in a CI job after the build step.

% ---------------------------------------------------------------------------
\section{Command line reference}
\label{sec:cli}

\begin{lstlisting}
sysinfo [OPTIONS]
\end{lstlisting}

\subsection{Options}

Options may appear in any order. Short options can be combined (\opt{-vr}) and may have their value attached
(\opt{-fjson}, \opt{-scpu,memory}). Long options accept \opt{--name value} and \opt{--name=value}.

\begin{longtable}{@{}p{4.6cm}p{10.4cm}@{}}
\toprule
Option & Description \\
\midrule
\endhead
\opt{-s, --section LIST} & Report only these sections. \texttt{LIST} is comma separated and the option may be
repeated. Names: \texttt{system cpu memory storage gpu network sensors environment build}; plus \texttt{core}
(system, cpu, memory, build) and \texttt{all}. Sections are always printed in the standard order, whatever the order
on the command line. \\
\opt{-x, --exclude LIST} & Leave these sections out; same syntax. Applied after \opt{--section}. \\
\opt{-l, --list} & Print the available sections (identifier, title, \texttt{(core)}) and exit. \\
\opt{-v, --verbose} & Add the long tail: the full CPU flag list, swappiness and overcommit policy, pseudo
file systems (\texttt{tmpfs}, \dots), disks such as loop devices, virtual network interfaces and link-local IPv6
addresses. \\
\midrule
\opt{-f, --format FORMAT} & \texttt{text} (default), \texttt{json} or \texttt{kv}. See Section~\ref{sec:formats}. \\
\opt{-j, --json} & Shortcut for \opt{--format json}. \\
\opt{-o, --output FILE} & Write to \texttt{FILE} (created or truncated) instead of standard output. Colour is then off. \\
\opt{-r, --redact} & Replace identifying values (hostname, user name, MAC and IP addresses) by
\texttt{<redacted>}; see Section~\ref{sec:privacy}. \\
\opt{--color[=WHEN]} & Colour the text format. \texttt{WHEN} is \texttt{auto} (default: only on a terminal, and
not if \texttt{NO\_COLOR} is set or \texttt{TERM=dumb}), \texttt{always} or \texttt{never}. Without a value:
\texttt{always}. \\
\opt{--no-color} & Same as \opt{--color=never}. \\
\midrule
\opt{-h, --help} & Show the usage summary and exit. \\
\opt{-V, --version} & Print \texttt{sysinfo 1.0.0} and exit. \\
\bottomrule
\end{longtable}

The built-in help text:

\lstinputlisting[style=out]{examples/help.txt}

\subsection{Exit status}

\begin{center}
\begin{tabular}{@{}cl@{}}
\toprule
Code & Meaning \\
\midrule
0 & success \\
1 & the output file could not be opened, or writing failed (disk full, \dots) \\
2 & usage error: unknown option or section, missing value, unexpected argument \\
\bottomrule
\end{tabular}
\end{center}

A piece of information that cannot be determined never makes the program fail: it is shown as
\texttt{unknown} or left out. Usage errors are reported on standard error.

\subsection{Environment}

\begin{description}[style=nextline]
  \item[\texttt{NO\_COLOR}] If set and non-empty, colour is not used in \texttt{auto} mode
        (see \url{https://no-color.org}).
  \item[\texttt{TERM}] \texttt{dumb} disables automatic colour.
\end{description}

\subsection{Examples}

\begin{lstlisting}
sysinfo -l                         # which sections exist?
sysinfo -s cpu                     # processor only
sysinfo -s core -x build           # system, cpu and memory
sysinfo -v -s cpu                  # processor, including all CPU flags
sysinfo -r > report.txt            # safe to paste into a bug report
sysinfo -j | jq '.sections.cpu.sys_cpu_model'
sysinfo -f kv -s memory | grep ram_total
\end{lstlisting}

% ---------------------------------------------------------------------------
\section{Output formats}
\label{sec:formats}

\subsection{Text}

Each section has a title; entries are listed as \emph{label}, \emph{value}. This is the format for reading.
Example (hostname and other identifying values redacted, abbreviated virtual machine):

\lstinputlisting[style=out]{examples/sample-text.txt}

\subsection{JSON}

One object with the program version, whether redaction was active, and a \texttt{sections} object that maps
each section identifier to an object of \emph{key} $\to$ \emph{value} strings. All values are strings; text is valid
UTF-8 (invalid bytes found in system files are replaced by U+FFFD).

\lstinputlisting[style=out]{examples/sample-json.json}

\subsection{Key/value}

One \texttt{key=value} per line, no quoting, in section order. Convenient for shell scripts and \cmd{grep}.

\lstinputlisting[style=out]{examples/sample-kv.txt}

The keys are the same in all formats; they are stable and always start with \texttt{sys\_}.

% ---------------------------------------------------------------------------
\section{What is reported}
\label{sec:sections}

All keys start with \texttt{sys\_}; the tables omit the prefix. Entries marked $\dagger$ are
optional: they are left out when the system does not provide them. Entries marked $^{*}$ are
identifying and are hidden by \opt{--redact}. Entries marked $^{v}$ appear only with \opt{--verbose}.

\subsection{system \textnormal{--- operating system and machine}}

\begin{longtable}{@{}p{5.7cm}p{10cm}@{}}
\toprule Key & Content \\ \midrule \endhead
\key{hostname}$^{*}$ & Host name. \\
\key{os}               & Distribution name and version from \cmd{/etc/os-release}, e.g.\ \texttt{Ubuntu 24.04.5 LTS}. \\
\key{kernel}           & Kernel name and release. \\
\key{arch}             & Machine architecture (\texttt{x86\_64}, \texttt{aarch64}, \dots). \\
\key{machine}$^\dagger$  & System vendor and product from DMI (or the device-tree model on ARM boards). \\
\key{board}$^\dagger$    & Mainboard vendor and name. \\
\key{firmware}$^\dagger$ & BIOS/UEFI vendor, version and date. \\
\key{virtualized}      & \texttt{yes} when the CPU reports a hypervisor or the DMI data names a virtual machine; otherwise
                         \texttt{no hypervisor flag}. \\
\key{container}$^\dagger$ & Docker, Podman, Kubernetes, LXC, \dots, when running inside one. \\
\key{wsl}$^\dagger$      & \texttt{yes} under Windows Subsystem for Linux. \\
\key{uptime}           & Time since boot. \\
\key{loadavg}          & Load average over 1, 5 and 15 minutes. \\
\bottomrule
\end{longtable}

\subsection{cpu \textnormal{--- processor}}

\begin{longtable}{@{}p{5.7cm}p{10cm}@{}}
\toprule Key & Content \\ \midrule \endhead
\key{cpu_model}              & Marketing name of the processor. \\
\key{cpu_vendor}$^\dagger$     & Vendor string (\texttt{GenuineIntel}, \texttt{AuthenticAMD}, \dots). \\
\key{cpu_signature}$^\dagger$  & x86 family / model / stepping. \\
\key{cpu_microcode}$^\dagger$  & Microcode revision. \\
\key{cpu_sockets}            & Number of processor packages. \\
\key{cpu_physical_cores}     & Number of physical cores (distinct hardware thread sibling groups). \\
\key{cpu_logical_cores}      & Number of online hardware threads. \\
\key{cpu_smt}$^\dagger$        & Simultaneous multithreading: \texttt{on}, \texttt{off} or \texttt{not supported}. \\
\key{cpu_numa_nodes}$^\dagger$ & Number of NUMA nodes. \\
\key{cpu_usable_cores}$^\dagger$ & Threads this process may run on, shown only when restricted (\cmd{taskset}, cpuset). \\
\key{cpu_cgroup_limit}$^\dagger$ & CPU quota of the container/cgroup, in CPUs. \\
\key{cpu_max_mhz}$^\dagger$\newline\key{cpu_min_mhz}$^\dagger$ & Frequency range allowed by the hardware, MHz (number only). \\
\key{cpu_cur_mhz}$^\dagger$    & Current frequency of cpu0, MHz (number only). \\
\key{cpu_governor}$^\dagger$   & Frequency governor of all CPUs; marked \texttt{(mixed)} if they differ. \\
\key{cpu_freq_driver}$^\dagger$ & Frequency scaling driver (\texttt{intel\_pstate}, \texttt{amd-pstate}, \dots). \\
\key{cpu_turbo}$^\dagger$      & Turbo/boost \texttt{enabled} or \texttt{disabled}. \\
\key{cpu_simd}$^\dagger$       & Notable instruction set extensions (SSE4.2, AVX2, AVX-512, AES-NI, NEON, SVE, \dots). \\
\key{cpu_flags}$^{v}$          & The complete flag list from \cmd{/proc/cpuinfo}. \\
\key{cpu_caches}$^\dagger$     & Cache sizes of cpu0, e.g.\ \texttt{L1d 48K, L1i 32K, L2 2048K, L3 266240K}. \\
\key{cpu_cache_line}$^\dagger$ & L1 data cache line size. \\
\bottomrule
\end{longtable}

\subsection{memory}

\begin{longtable}{@{}p{5.7cm}p{10cm}@{}}
\toprule Key & Content \\ \midrule \endhead
\key{ram_total}                & Installed RAM as seen by the kernel (GiB, one decimal). \\
\key{ram_available}$^\dagger$    & RAM available to new programs without swapping. \\
\key{ram_cgroup_limit}$^\dagger$ & Memory limit of the container/cgroup, when lower than the physical RAM. \\
\key{swap}$^\dagger$             & Swap size and use, or \texttt{none}. \\
\key{thp}$^\dagger$              & Transparent huge pages mode; the active one is in brackets. \\
\key{hugepages}$^\dagger$        & Huge page size and the number of reserved pages. \\
\key{page_size}$^\dagger$        & Virtual memory page size. \\
\key{overcommit}$^{v}$\newline\key{swappiness}$^{v}$ & Kernel tunables \cmd{vm.overcommit\_memory} and \cmd{vm.swappiness}. \\
\bottomrule
\end{longtable}

\subsection{storage}

\begin{longtable}{@{}p{5.7cm}p{10cm}@{}}
\toprule Key & Content \\ \midrule \endhead
\key{disk_}\textit{name} & One per block device: size, kind (\texttt{NVMe SSD}, \texttt{SSD}, \texttt{HDD},
\texttt{virtual disk}, \texttt{MMC/SD}) and model. Loop, RAM, optical and device-mapper devices only with \opt{--verbose}. \\
\key{fs_}\textit{mount}  & One per mounted file system: type, source device, size and percentage used. By default only
file systems backed by storage (ext4, xfs, btrfs, zfs, vfat, ntfs, NFS, \dots) and the root overlay of a container; bind
mounts of the same device are listed once. The key is derived from the mount point (\texttt{fs\_root} for \texttt{/}). \\
\bottomrule
\end{longtable}

\subsection{gpu \textnormal{--- graphics}}

\begin{longtable}{@{}p{5.7cm}p{10cm}@{}}
\toprule Key & Content \\ \midrule \endhead
\key{gpu_}\textit{n}              & One per display controller on the PCI bus: vendor and model (looked up in the system's
\cmd{pci.ids} database, or the NVIDIA driver's model name), kernel driver, and video memory where the driver exposes it
(AMD). Without PCI graphics, DRM devices of system-on-chip boards are listed; otherwise \texttt{none detected}. \\
\key{gpu_nvidia_driver}$^\dagger$ & Version of the proprietary NVIDIA kernel module. \\
\bottomrule
\end{longtable}

\subsection{network}

\begin{longtable}{@{}p{5.7cm}p{10cm}@{}}
\toprule Key & Content \\ \midrule \endhead
\key{net_}\textit{interface}$^{*}$ & State, wireless marker, link speed, MTU, driver, MAC address and IPv4/IPv6 addresses with prefix
length. With \opt{--redact} the MAC and the addresses are removed, the rest is kept. The loopback device, Docker/KVM bridges and
\texttt{veth} pairs are hidden unless \opt{--verbose} is given. \\
\bottomrule
\end{longtable}

\subsection{sensors \textnormal{--- power and temperature}}

\begin{longtable}{@{}p{5.7cm}p{10cm}@{}}
\toprule Key & Content \\ \midrule \endhead
\key{battery_}\textit{name}   & Charge, status, technology, health (full charge relative to design capacity), cycle count, model. \\
\key{ac_}\textit{name}        & Mains adapter connected or not. \\
\key{temp_}\textit{chip}      & Temperature of each hardware monitor chip (\texttt{coretemp}, \texttt{k10temp}, \texttt{nvme}, \dots); thermal
zones are used when no chip is present. \\
\key{sensors}               & Shown instead when nothing is exposed (typical for virtual machines and containers). \\
\bottomrule
\end{longtable}

\subsection{environment}

\begin{longtable}{@{}p{5.7cm}p{10cm}@{}}
\toprule Key & Content \\ \midrule \endhead
\key{user}$^{*}$\newline\key{is_root} & Effective user name and whether it is the superuser. \\
\key{shell}\newline\key{lang}\newline\key{term} & \texttt{\$SHELL}, \texttt{\$LANG}, \texttt{\$TERM}. \\
\key{session_type}\newline\key{desktop}  & \texttt{\$XDG\_SESSION\_TYPE} (x11/wayland) and \texttt{\$XDG\_CURRENT\_DESKTOP}. \\
\key{timezone}                    & From \cmd{/etc/timezone} or the \cmd{/etc/localtime} link. \\
\bottomrule
\end{longtable}
Entries whose variable is unset are left out.

\subsection{build \textnormal{--- how this program was built}}

\begin{longtable}{@{}p{5.7cm}p{10cm}@{}}
\toprule Key & Content \\ \midrule \endhead
\key{compiler}\newline\key{stdlib}    & Compiler and C++ standard library with versions. \\
\key{cxx_std}                   & C++ language level actually compiled, e.g.\ \texttt{C++20 (202002L)}. \\
\key{build_type}\newline\key{cxx_flags} & CMake build type and the compiler flags in effect. \\
\key{pointer_bits}\newline\key{endianness} & Word size and byte order of this binary. \\
\key{assertions}                & Whether \cmd{assert} is active (\texttt{NDEBUG}). \\
\key{libc}$^\dagger$              & C library version (glibc). \\
\key{sysinfo_version}           & Version of the collector. \\
\bottomrule
\end{longtable}

% ---------------------------------------------------------------------------
\section{Recipes}

\subsection{Checking a machine before benchmarking}

Timing results are only comparable when the machine is in a known state. Look at:

\begin{lstlisting}
sysinfo -s cpu,memory,sensors | grep -Ei 'governor|turbo|smt|frequency|limit|cur'
\end{lstlisting}

\begin{itemize}
  \item \texttt{cpu\_governor} should be \texttt{performance} and not \texttt{(mixed)}; \texttt{powersave} lets the clock drift.
  \item \texttt{cpu\_turbo} and \texttt{cpu\_smt} should be what you intend; decide once and keep it.
  \item \texttt{cpu\_cgroup\_limit}, \texttt{ram\_cgroup\_limit} and \texttt{cpu\_usable\_cores} reveal a container or \cmd{taskset} that
        restricts the program.
  \item \texttt{virtualized} \texttt{yes} means timings are subject to the host's scheduling.
  \item \texttt{build\_type} must be \texttt{Release} or similar; \texttt{assertions} should be \texttt{disabled}.
\end{itemize}

\subsection{Recording the machine next to results}

\begin{lstlisting}
sysinfo -j -o results/machine.json
sysinfo -f kv -s core -r > results/machine.kv
\end{lstlisting}

\subsection{Comparing two machines}

\begin{lstlisting}
sysinfo -f kv -r -x environment,network > a.kv      # on machine A
sysinfo -f kv -r -x environment,network > b.kv      # on machine B
diff a.kv b.kv
\end{lstlisting}

\subsection{Using values in a script}

\begin{lstlisting}
cores=$(sysinfo -f kv -s cpu | sed -n 's/^sys_cpu_physical_cores=//p')
make -j"$cores"
\end{lstlisting}

\subsection{Reporting a bug}

Attach the output of \cmd{sysinfo -v -r}; the hostname, user name, MAC and IP addresses are removed.

% ---------------------------------------------------------------------------
\section{Using the library}
\label{sec:library}

The collector is a separate static library, \texttt{sysinfo::sysinfo}, that the command line tool is built on. Use it from
another CMake project:

\begin{lstlisting}[style=cmake]
add_subdirectory(external/sysinfo)                 # or FetchContent
target_link_libraries(my_target PRIVATE sysinfo::sysinfo)
\end{lstlisting}

\begin{lstlisting}[style=cpp]
#include <sysinfo.h>

// structured: sections with titled entries
for (const sysinfo::Section& s : sysinfo::collectAll()) {
  std::cout << s.title << '\n';
  for (const sysinfo::Entry& e : s.entries)
    std::cout << "  " << e.label << ": " << e.display(/*redact=*/false) << '\n';
}

// flat: the "sys_*" pairs of the core sections
sysinfo::KeyValues kv = sysinfo::collect();
// e.g. with Google Benchmark:
for (const auto& [key, value] : kv) benchmark::AddCustomContext(key, value);
\end{lstlisting}

\begin{center}
\begin{tabular}{@{}p{6.2cm}p{8.6cm}@{}}
\toprule
Function & Purpose \\
\midrule
\texttt{version()}                    & Library version string. \\
\texttt{sections()}                   & Identifier, title and ``core'' flag of every section. \\
\texttt{collectSection(id, verbose)}  & One section; an unknown identifier gives an empty section. \\
\texttt{collectAll(verbose)}          & All sections in display order. \\
\texttt{collect()}                    & Flat list for the core sections (system, cpu, memory, build). \\
\bottomrule
\end{tabular}
\end{center}

None of the functions throws; anything that fails is reported as \texttt{unknown}, left out, or (for a whole section)
shown as a single \texttt{could not be collected} entry.

\textbf{Build information.} CMake generates \cmd{build\_info.h} with the build type and compiler flags.

% ---------------------------------------------------------------------------
\section{How the information is obtained}

\begin{center}
\small
\begin{tabular}{@{}p{3.0cm}p{12cm}@{}}
\toprule
Section & Sources \\
\midrule
system      & \path{/etc/os-release}, \texttt{uname}, \path{/sys/devices/virtual/dmi/id/*}, \path{/proc/uptime}, \path{/proc/loadavg}, \path{/proc/1/cgroup} \\
cpu         & \path{/proc/cpuinfo}, \path{/sys/devices/system/cpu/*} (topology, cpufreq, cache, smt, intel\_pstate), \path{/sys/devices/system/node},
              \texttt{sched\_getaffinity}, \path{/sys/fs/cgroup/cpu.max} \\
memory      & \path{/proc/meminfo}, \path{/sys/kernel/mm/transparent_hugepage}, \path{/sys/fs/cgroup/memory.max}, \texttt{sysconf} \\
storage     & \path{/sys/block/*}, \path{/proc/self/mounts}, \texttt{statvfs} \\
gpu         & \path{/sys/bus/pci/devices/*}, \texttt{pci.ids} (\path{/usr/share/hwdata}, \path{/usr/share/misc}), \path{/proc/driver/nvidia}, \path{/sys/class/drm} \\
network     & \path{/sys/class/net/*}, \texttt{getifaddrs} \\
sensors     & \path{/sys/class/power_supply}, \path{/sys/class/hwmon}, \path{/sys/class/thermal} \\
environment & environment variables, \texttt{getpwuid}, \path{/etc/timezone} \\
\bottomrule
\end{tabular}
\end{center}

\subsection{Limits to be aware of}

\begin{itemize}
  \item \textbf{No root, no external tools.} Information that needs administrator rights is not reported:
        memory module type and speed (\cmd{dmidecode}), serial numbers, SMART data, GPU load and clocks.
  \item \textbf{Virtual machines and containers} show what the guest sees. The CPU model may be generic, frequency and
        temperature files are often missing, and the cache sizes may be those of the host.
  \item \textbf{Caches} are read for cpu0 only. On hybrid processors (performance and efficiency cores) other cores can differ.
  \item \textbf{Current frequency} is a momentary reading of cpu0.
  \item \textbf{cgroup limits} are read for the program's own cgroup and the root; limits set on intermediate parent groups
        are not followed.
  \item \textbf{GPU names} depend on an up-to-date \cmd{pci.ids} (package \texttt{hwdata} or \texttt{pci.ids}); otherwise the
        numeric PCI identifier is shown.
  \item \textbf{Other platforms.} On macOS (via \cmd{sysctl}) and Windows (Win32 API) only the system, cpu, memory and build
        sections are filled, with fewer entries; the remaining sections say \texttt{not available on this platform}. These
        ports compile conditionally and have not been tested as thoroughly as Linux.
\end{itemize}

% ---------------------------------------------------------------------------
\section{Privacy}
\label{sec:privacy}

\prog{} only reads local files and never sends anything anywhere. The output can nevertheless identify the machine or
its user: the hostname, the user name, the MAC and IP addresses, the board and firmware version.
\opt{--redact} replaces \texttt{hostname} and \texttt{user} by \texttt{<redacted>} and removes the MAC and IP addresses from the
network entries. It does \emph{not} remove other details such as the exact machine model, the time zone or the kernel version; use
\opt{-x network,environment,system} to leave whole sections out when in doubt.

% ---------------------------------------------------------------------------
\section{Troubleshooting}
\label{sec:trouble}

\begin{description}[style=nextline,leftmargin=0.6cm]
  \item[CMake: ``Target requires the language dialect CXX20''] The compiler is too old. Use GCC 11+ or Clang 14+
        (\cmd{-DCMAKE\_CXX\_COMPILER=g++-13}).
  \item[Many entries say \texttt{unknown}] The file under \cmd{/proc} or \cmd{/sys} does not exist or is not readable. This is normal in
        containers, virtual machines and restricted sandboxes. Compare with \cmd{sysinfo -v}.
  \item[No frequency, governor or turbo lines] The kernel exposes no \texttt{cpufreq} interface (virtual machine, or the driver is
        not loaded).
  \item[\texttt{cpu\_physical\_cores} equals \texttt{cpu\_logical\_cores}] The CPU has no SMT, SMT is disabled, or the hypervisor presents
        each thread as a core.
  \item[Garbage escape codes in a file or pipe] Colour is only enabled automatically on a terminal; add \opt{--no-color} or set \texttt{NO\_COLOR}
        when a wrapper forces it.
  \item[GPU shows only hexadecimal numbers] Install the PCI ID database: \cmd{sudo apt install hwdata} (or \cmd{pci.ids}).
  \item[Exit status 2] Read the message on standard error; \opt{--help} lists the valid options and \opt{--list} the sections.
\end{description}

\end{document}
